% Options for packages loaded elsewhere
\PassOptionsToPackage{unicode}{hyperref}
\PassOptionsToPackage{hyphens}{url}
\PassOptionsToPackage{dvipsnames,svgnames,x11names}{xcolor}
\documentclass[
  11pt,
  a4paper,
]{article}
\usepackage{xcolor}
\usepackage[margin=2.6cm]{geometry}
\usepackage{amsmath,amssymb}
\setcounter{secnumdepth}{5}
\usepackage{iftex}
\ifPDFTeX
  \usepackage[T1]{fontenc}
  \usepackage[utf8]{inputenc}
  \usepackage{textcomp} % provide euro and other symbols
\else % if luatex or xetex
  \usepackage{unicode-math} % this also loads fontspec
  \defaultfontfeatures{Scale=MatchLowercase}
  \defaultfontfeatures[\rmfamily]{Ligatures=TeX,Scale=1}
\fi
\usepackage{lmodern}
\ifPDFTeX\else
  % xetex/luatex font selection
  \setmainfont[]{Libertinus Serif}
  \setmonofont[Scale=0.85]{Latin Modern Mono}
  \setmathfont[]{Latin Modern Math}
\fi
% Use upquote if available, for straight quotes in verbatim environments
\IfFileExists{upquote.sty}{\usepackage{upquote}}{}
\IfFileExists{microtype.sty}{% use microtype if available
  \usepackage[]{microtype}
  \UseMicrotypeSet[protrusion]{basicmath} % disable protrusion for tt fonts
}{}
\makeatletter
\@ifundefined{KOMAClassName}{% if non-KOMA class
  \IfFileExists{parskip.sty}{%
    \usepackage{parskip}
  }{% else
    \setlength{\parindent}{0pt}
    \setlength{\parskip}{6pt plus 2pt minus 1pt}}
}{% if KOMA class
  \KOMAoptions{parskip=half}}
\makeatother
\usepackage{longtable,booktabs,array}
\newcounter{none} % for unnumbered tables
\usepackage{calc} % for calculating minipage widths
% Correct order of tables after \paragraph or \subparagraph
\usepackage{etoolbox}
\makeatletter
\patchcmd\longtable{\par}{\if@noskipsec\mbox{}\fi\par}{}{}
\makeatother
% Allow footnotes in longtable head/foot
\IfFileExists{footnotehyper.sty}{\usepackage{footnotehyper}}{\usepackage{footnote}}
\makesavenoteenv{longtable}
\setlength{\emergencystretch}{3em} % prevent overfull lines
\providecommand{\tightlist}{%
  \setlength{\itemsep}{0pt}\setlength{\parskip}{0pt}}

\providecommand{\E}{\mathbb{E}}
\providecommand{\N}{\mathbb{N}}
\providecommand{\Z}{\mathbb{Z}}
\providecommand{\R}{\mathbb{R}}
\providecommand{\eps}{\varepsilon}
\providecommand{\ceil}[1]{\left\lceil #1\right\rceil}
\providecommand{\floor}[1]{\left\lfloor #1\right\rfloor}
\AtBeginDocument{\let\setminus\smallsetminus}
\usepackage[most]{tcolorbox}
\newtcolorbox{warning}{colback=red!5,colframe=red!65!black,boxrule=0.8pt,arc=2pt,left=6pt,right=6pt,top=5pt,bottom=5pt}
\usepackage{etoolbox}
\AtBeginEnvironment{longtable}{\small}
\setlength{\emergencystretch}{3em}
\usepackage{fancyhdr}
\pagestyle{fancy}
\fancyhf{}
\fancyhead[L]{\small\itshape Candidate result 2603.02786\_\_01 --- GPT-6 Astra, awaiting review}
\fancyhead[R]{\small\thepage}
\renewcommand{\headrulewidth}{0.2pt}
\usepackage{bookmark}
\IfFileExists{xurl.sty}{\usepackage{xurl}}{} % add URL line breaks if available
\urlstyle{same}
\hypersetup{
  pdftitle={A candidate proof of Conjecture 4 in Packing arithmetic progressions},
  pdfauthor={github.com/graph-theory-AI --- machine-generated candidate writeup},
  colorlinks=true,
  linkcolor={blue!50!black},
  filecolor={Maroon},
  citecolor={Blue},
  urlcolor={blue!50!black},
  pdfcreator={LaTeX via pandoc}}

\title{A candidate proof of Conjecture 4 in Packing arithmetic
progressions}
\author{\href{https://github.com/graph-theory-AI}{github.com/graph-theory-AI}
--- machine-generated candidate writeup}
\date{17 September 2026}

\begin{document}
\maketitle
\begin{abstract}
The conjectured asymptotic follows from an exact packing lemma for
arithmetic-progression blocks of differences and the dense-subset
Green--Tao theorem. This candidate proof was generated by GPT-6 Astra in
a single model pass for the Graph Theory AI project. It was selected
from the campaign because the model marked it
\texttt{would\_publish:\ true}; it has not been checked by the
adversarial referee, a human mathematician, or a novelty search.
\end{abstract}

\section{The problem and the claimed
result}\label{the-problem-and-the-claimed-result}

This document concerns
\href{https://graph-theory-ai.github.io/graph-conjectures/arxiv/2603.02786__01/}{Conjecture
4} from \emph{Packing arithmetic progressions}.

\subsection{Problem statement}\label{problem-statement}

Conjecture.
\(M_{\mathbb{P}(n)}(n)\sim\frac{1}{6}\cdot\frac{n^{3}}{\ln n}\).

\subsection{Catalog context}\label{catalog-context}

Closely related to Conjecture 3. The lower bounds in Conjectures 3 and 4
follow from the proofs of the lower bounds in Theorems 1 and 5,
respectively. The upper bounds remain open.

\subsection{Claimed result}\label{claimed-result}

\textbf{Model verdict: proved} (confidence: high).

The conjectured asymptotic follows from an exact packing lemma for
arithmetic-progression blocks of differences and the dense-subset
Green--Tao theorem.

\emph{The model's own caveat:} Uses the established relative-density
Green--Tao theorem; novelty has not been checked against the full source
paper.

\section{Model writeup}\label{model-writeup}

The text below is the model output verbatim, apart from moving its
machine-readable verdict block into the summary above and shifting
Markdown heading levels for this document. Mechanical TeX defects and
equation tags were normalized where needed for compilation.

\subsection{Proof of the conjecture}\label{proof-of-the-conjecture}

The restriction in the previous attempt---that nonconsecutive convex
hulls must be disjoint---can be bypassed. For suitable
arithmetic-progression blocks of prime differences, one can coordinate
\textbf{every pair} of progressions, even when many convex hulls
overlap. The dense-subset Green--Tao theorem supplies a partition of
almost all primes into such blocks.

All shifts below are integers. An integer interval of length \(L\)
contains \(L\) consecutive integers; using diameter instead changes the
packing parameter by one.

Write \[
A_d(b)=\{b,b+d,\ldots,b+(n-1)d\},
\qquad
S(n)=\sum_{p\le n}p(n-p).
\]

We prove the slightly more informative statement \[
\boxed{\;
M_{\mathbb P(n)}(n)
=
S(n)+o\!\left(\frac{n^3}{\log n}\right).
\;}
\] The prime number theorem then gives the conjectured constant \(1/6\).

\subsubsection{1. An exact pairwise avoidance
criterion}\label{an-exact-pairwise-avoidance-criterion}

\paragraph{Lemma 1}\label{lemma-1}

Let \(p,q\) be coprime positive integers. Suppose \[
b'-b=Ap+Bq
\] for integers \(A,B\) satisfying \[
A\ge n-q,\qquad B\ge1.
\] Then \(A_p(b)\) and \(A_q(b')\) are disjoint.

\paragraph{Proof}\label{proof}

An intersection would give \[
b+up=b'+vq
\] for some \(0\le u,v<n\), and hence \[
(u-A)p=(v+B)q.
\] Since the right-hand side is positive and \(\gcd(p,q)=1\), the
integer \(u-A\) is a positive multiple of \(q\). Consequently, \[
u\ge A+q\ge n,
\] a contradiction. \(\square\)

The next lemma arranges for this criterion to hold simultaneously for
every pair in a block.

\subsubsection{2. Packing an arithmetic-progression block of
differences}\label{packing-an-arithmetic-progression-block-of-differences}

\paragraph{Lemma 2}\label{lemma-2}

Fix \(k\ge2\), and put \[
Q=\operatorname{lcm}(1,2,\ldots,k-1).
\] Suppose \[
d_i=a+ih,\qquad 0\le i<k,
\] are pairwise coprime positive integers satisfying \[
d_0<d_1<\cdots<d_{k-1}\le n,
\qquad Q\mid h.
\] Then their \(n\)-term progressions can be packed into an integer
interval of length at most \[
\boxed{\;
\sum_{i=0}^{k-1}d_i(n-d_i)+n^2+2Qkn.
\;}
\]

\paragraph{Proof}\label{proof-1}

Choose the integer \(n_*\) satisfying \[
n\le n_*\le n+Q-1,
\qquad n_*\equiv a\pmod Q.
\] Thus every \(n_*-d_i\) is a nonnegative multiple of \(Q\).

Set \(b_0=0\), and recursively define \[
b_{i+1}
=
b_i+d_i(n_*-d_{i+1})+Qd_{i+1}
\qquad(0\le i<k-1).
\] We claim that all sets \(A_{d_i}(b_i)\) are pairwise disjoint.

Fix \(i<j\), and write \[
m=j-i,\qquad p=d_i,\qquad q=d_j.
\] For \(0\le t<m\), the arithmetic-progression structure gives \[
d_{i+t}=\frac{m-t}{m}p+\frac{t}{m}q
\] and \[
d_{i+t+1}
=\frac{m-t-1}{m}p+\frac{t+1}{m}q.
\] Therefore \begin{equation*}
\begin{aligned}
d_{i+t}(n_*-d_{i+t+1})
&=
\frac{(m-t)(n_*-d_{i+t+1})}{m}\,p\\
&\quad+
\frac{t(n_*-d_{i+t+1})}{m}\,q,
\end{aligned}
\end{equation*} while \[
Qd_{i+t+1}
=
\frac{Q(m-t-1)}m\,p
+
\frac{Q(t+1)}m\,q.
\]

\textbf{Every coefficient displayed here is a nonnegative integer.}
Indeed, \(m\le k-1\), so \(m\mid Q\), and \(Q\mid n_*-d_{i+t+1}\).

Summing these identities over \(t=0,\ldots,m-1\), we obtain \[
b_j-b_i=Ap+Bq
\] with integer coefficients. The first term of the first sum
contributes \((n_*-d_{i+1})p\), so \[
A\ge n_*-d_{i+1}\ge n-q.
\] The last term of the second sum contributes \(Qq\), so \[
B\ge Q\ge1.
\] Lemma 1 now proves disjointness for this pair. Since \(i<j\) was
arbitrary, all pairs are covered.

It remains to bound the containing interval. The \(b_i\) increase
strictly, and the \(d_i\) increase, so the largest endpoint belongs to
the last progression. Its containing interval has length \[
\Lambda
=
\sum_{i=0}^{k-2}
\bigl[d_i(n_*-d_{i+1})+Qd_{i+1}\bigr]
+(n-1)d_{k-1}+1.
\] For \(i<k-1\), using \(d_{i+1}\ge d_i\), \[
d_i(n_*-d_{i+1})
\le d_i(n-d_i)+(n_*-n)d_i.
\] Also, \[
(n-1)d_{k-1}+1
=
d_{k-1}(n-d_{k-1})
+d_{k-1}^2-d_{k-1}+1
\le d_{k-1}(n-d_{k-1})+n^2.
\] Because \(n_*-n<Q\) and all \(d_i\le n\), these inequalities give \[
\Lambda
\le
\sum_{i=0}^{k-1}d_i(n-d_i)+n^2+2Qkn.
\] This proves the lemma. \(\square\)

The crucial feature is the cost: a block of \(k\) differences incurs
only one \(n^2\)-sized boundary overhead, plus \(O_k(n)\). No separation
condition on nonconsecutive convex hulls is imposed.

\subsubsection{3. Almost all primes can be partitioned into suitable
blocks}\label{almost-all-primes-can-be-partitioned-into-suitable-blocks}

The non-elementary input is the following established form of the
Green--Tao theorem.

\paragraph{Dense-subset Green--Tao
theorem}\label{dense-subset-greentao-theorem}

For every fixed \(k\ge3\) and \(\delta>0\), all sufficiently large \(N\)
have the following property: every set \[
E\subseteq\mathbb P(N),
\qquad |E|\ge\delta\,\pi(N),
\] contains a nonconstant \(k\)-term arithmetic progression.

The dense-subset form matters here: the mere existence of arbitrarily
long prime arithmetic progressions would not suffice.

\paragraph{Lemma 3}\label{lemma-3}

For any fixed positive integers \(k\ge3\) and \(Q\), all but \[
o_{k,Q}(\pi(n))
\] primes up to \(n\) can be partitioned into \(k\)-term arithmetic
progressions whose common differences are divisible by \(Q\).

\paragraph{Proof}\label{proof-2}

Take a maximal collection of disjoint \(k\)-term arithmetic progressions
of primes up to \(n\), each with common difference divisible by \(Q\).
Let \(R_n\) be the unused primes.

If \(|R_n|\ge\delta\pi(n)\), some residue class modulo \(Q\) contains at
least \[
\frac{\delta}{Q}\pi(n)
\] members of \(R_n\). For sufficiently large \(n\), the dense-subset
Green--Tao theorem supplies a \(k\)-term arithmetic progression within
that residue class. Its common difference is divisible by \(Q\),
contradicting maximality.

This works for every fixed \(\delta>0\), proving the claim. \(\square\)

\subsubsection{4. The asymptotically optimal upper
bound}\label{the-asymptotically-optimal-upper-bound}

Fix \(k\ge3\), and let \[
Q=\operatorname{lcm}(1,\ldots,k-1).
\] Apply Lemma 3. Let \(t_n\) be the number of resulting blocks and
\(u_n\) the number of unused primes. Then \[
kt_n+u_n=\pi(n),
\qquad
u_n=o_k(\pi(n)).
\]

Each block consists of distinct primes, hence its differences are
pairwise coprime. Lemma 2 applies.

Pack the blocks into consecutive, disjoint integer intervals. Pack each
unused prime progression into its own interval. For a singleton prime
\(p\le n\), the required length satisfies \[
(n-1)p+1
=
p(n-p)+(p^2-p+1)
\le p(n-p)+n^2.
\] Adding all block and singleton costs gives \begin{equation*}
\begin{aligned}
M_{\mathbb P(n)}(n)
&\le S(n)+(t_n+u_n)n^2+2Qkn\,t_n\\
&\le S(n)+\frac{n^2\pi(n)}k
       +u_n n^2+2Qn\pi(n).
\end{aligned}
\end{equation*} Here \(k\) and \(Q\) are fixed. Consequently, \[
\limsup_{n\to\infty}
\frac{M_{\mathbb P(n)}(n)-S(n)}{n^2\pi(n)}
\le \frac1k.
\] Since this holds for every fixed \(k\ge3\), \[
\boxed{\;
\limsup_{n\to\infty}
\frac{M_{\mathbb P(n)}(n)-S(n)}{n^2\pi(n)}
\le0.
\;}
\]

The order of limits is important: first \(n\to\infty\) with \(k\) fixed,
and only then \(k\to\infty\). Thus the rapid growth of \(Q\), and the
lack of a useful quantitative Green--Tao bound, cause no difficulty.

\subsubsection{5. Matching lower bound}\label{matching-lower-bound}

For completeness, here is an independent proof of the required lower
bound.

Consider any packing, and restrict it to the primes \(p<n\). Order these
progressions by their first terms: \[
a_1<a_2<\cdots<a_s,
\] with corresponding differences \(p_1,\ldots,p_s\).

For consecutive progressions, put \(p=p_i\), \(q=p_{i+1}\). If \[
a_{i+1}-a_i\le (n-1)p-pq,
\] then the integer interval \[
[a_{i+1},\,a_{i+1}+pq-1]
\] lies inside both convex hulls. For the second hull, this uses \[
(n-1)q\ge pq,
\] since \(p\le n-1\). The Chinese remainder theorem supplies a point in
this interval congruent to \(a_i\pmod p\) and to \(a_{i+1}\pmod q\).
Such a point belongs to both progressions, a contradiction.

Hence \[
a_{i+1}-a_i>(n-1)p_i-p_ip_{i+1}.
\] Summing and adding the span of the last progression yields the
weaker, sufficient inequality \[
M_{\mathbb P(n)}(n)
\ge
(n-1)\sum_{i=1}^s p_i
-\sum_{i=1}^{s-1}p_ip_{i+1}.
\] Using \(2xy\le x^2+y^2\), \[
\sum_{i=1}^{s-1}p_ip_{i+1}
\le \sum_{i=1}^s p_i^2.
\] It follows that \[
M_{\mathbb P(n)}(n)
\ge
\sum_{p<n}p(n-p)-\sum_{p<n}p.
\] If \(n\) itself is prime, its contribution to \(S(n)\) is zero.
Therefore \[
M_{\mathbb P(n)}(n)\ge S(n)-n\pi(n).
\] In particular, \[
\liminf_{n\to\infty}
\frac{M_{\mathbb P(n)}(n)-S(n)}{n^2\pi(n)}
\ge0.
\]

Together with the upper bound, this proves \[
M_{\mathbb P(n)}(n)
=
S(n)+o\bigl(n^2\pi(n)\bigr).
\]

\subsubsection{6. Evaluating the main
term}\label{evaluating-the-main-term}

By the prime number theorem and partial summation, \[
\sum_{p\le n}p
=
\left(\frac12+o(1)\right)\frac{n^2}{\log n},
\qquad
\sum_{p\le n}p^2
=
\left(\frac13+o(1)\right)\frac{n^3}{\log n}.
\] Thus \begin{equation*}
\begin{aligned}
S(n)
&=n\sum_{p\le n}p-\sum_{p\le n}p^2\\
&=
\left(\frac16+o(1)\right)\frac{n^3}{\log n}.
\end{aligned}
\end{equation*} Also \(n^2\pi(n)\sim n^3/\log n\). Consequently, \[
\boxed{\displaystyle
M_{\mathbb P(n)}(n)\sim \frac16\,\frac{n^3}{\log n}.
}
\]

All collision cases, including nonconsecutive progressions and primes
above \(n/2\), are covered by Lemma 2. The proof uses no unproved
prime-pattern hypothesis; its only deep input is the established
dense-subset Green--Tao theorem.

\appendix

\section{Provenance and review
status}\label{provenance-and-review-status}

\begin{warning}

\textbf{UNREFEREED MODEL OUTPUT.} This document was selected solely
because GPT-6 Astra labelled its own result
\texttt{would\_publish:\ true} and returned \texttt{proved} or
\texttt{disproved}. It has not passed the adversarial LLM referee used
for the earlier Sol campaign, has not been checked by a human
mathematician, and has not been checked for novelty. Treat every
mathematical and bibliographic claim as unverified.

\end{warning}

{\def\LTcaptype{none} % do not increment counter
\begin{longtable}[]{@{}
  >{\raggedright\arraybackslash}p{(\linewidth - 2\tabcolsep) * \real{0.5000}}
  >{\raggedright\arraybackslash}p{(\linewidth - 2\tabcolsep) * \real{0.5000}}@{}}
\toprule\noalign{}
\endhead
\bottomrule\noalign{}
\endlastfoot
Catalog id & \texttt{2603.02786\_\_01} \\
Catalog entry &
\href{https://graph-theory-ai.github.io/graph-conjectures/arxiv/2603.02786__01/}{Conjecture
4} \\
Source corpus & arxiv \\
Campaign leg & selected second attempts (\texttt{attacks\_retry}) \\
Source paper / entry & Packing arithmetic progressions \\
Model verdict & \textbf{proved} (confidence: high) \\
Selection criterion & \texttt{would\_publish:\ true} and verdict
\texttt{proved} or \texttt{disproved} \\
Model & \texttt{gpt-6-astra}, reasoning effort \texttt{max},
\texttt{mode=pro}, flex service tier \\
Original artifact &
\texttt{attacks\_retry/2603.02786\_\_01/output.md} \\
Independent review & \textbf{None} \\
arXiv & \href{https://arxiv.org/abs/2603.02786}{arXiv:2603.02786} \\
Previous attempt & \texttt{attacks/2603.02786\_\_01}: gpt-5.6-sol
verdict \texttt{partial} \\
Model's caveats & Uses the established relative-density Green--Tao
theorem; novelty has not been checked against the full source paper. \\
\end{longtable}
}

The generated Markdown and LaTeX sources for this PDF are stored under
\texttt{to\_review\_astra/src/2603.02786\_\_01/}.

\end{document}
